\documentclass[40pt]{article}

% Language setting
% Replace `english' with e.g. `spanish' to change the document language
\usepackage[english]{babel}

% Set page size and margins
% Replace `letterpaper' with `a4paper' for UK/EU standard size
\usepackage[letterpaper,top=2cm,bottom=2cm,left=3cm,right=3cm,marginparwidth=1.75cm]{geometry}

% Useful packages
\usepackage{amsmath}
\usepackage{amsthm}
\usepackage{amssymb}
\usepackage{graphicx}
\usepackage[colorlinks=true, allcolors=blue]{hyperref}
\usepackage{xcolor}
\usepackage{quiver}

\setlength{\parindent}{0cm}
\newtheorem{definition}{Definition}[section]
\newtheorem{theorem}{Theorem}[section]
\newcommand{\bb}[1]{\mathbb{#1}}
\newcommand{\mcal}[1]{\mathcal{#1}}
\newcommand{\R}{\mathbb{R}}
\newcommand{\Proj}{\mathbb{P}}
\newcommand{\N}{\mathbb{N}}
\newcommand{\C}{\mathbb{C}}
\newcommand{\A}{\mathbb{A}}
\newcommand{\Z}{\mathbb{Z}}
\newcommand{\Q}{\mathbb{Q}}
\newcommand{\I}{\mathcal{I}}
\newcommand{\V}{\mathcal{V}}
\newcommand{\la}{\langle}
\newcommand{\ra}{\rangle}
\newcommand{\im}{\text{im}\,}
\DeclareMathOperator{\id}{\text{id}}
\DeclareMathOperator{\rank}{\text{rk}}
\DeclareMathOperator{\Frac}{\text{Frac}}
\DeclareMathOperator{\trdeg}{\text{trdeg}}
\renewcommand{\qedsymbol}{$\aleph$}
\newcommand{\note}[1]{\textcolor{red}{\emph{#1}}}
%\newtheorem{theorem}{Theorem}[section]
\newtheorem{lemma}[theorem]{Lemma}
\newtheorem{remark}[theorem]{Remark}
%\newtheorem{definition}[theorem]{Definition}
\title{MATH 215A : Algebraic Topology}
\date{August 2026}


\begin{document}
	\maketitle
	\section{Homotopy Theory}
	Begin with a slew of definitions and attemps for each to furnish examples
	\begin{definition}
		We say that two maps $f_0, f_1 : X \rightarrow Y$ are homotopic, denoted $f_0 \sim f_1$, if there exists a map $F: X \times [0,1] \rightarrow Y$ such that the restriction $f_t = F|_{X \times \{t\}}$ corresponds to $f_0$ when $t = 0$ and to $f_1$ when $t = 1$. 
	\end{definition}
	Intuitively we are saying that two paths are homotopic to one another if they can be continuously deformed into one another. Extending this example we can determine when two spaces are homotopic to one another. 
	\begin{definition}
		$X$ is homotopic to $Y$ if there exists continuous maps $f: X \rightarrow Y : g$ such that $f \circ g \sim \id_X$ and $g \circ f \sim \id_Y$. 
	\end{definition}
	Note that if instead of homotopic we get equivalent then this is the meaning of being homeomorphic (bijective, continuous maps). In a sense we are presenting a weaker form of equivalent of two spaces. Some examples of homotopic spaces:
	\begin{enumerate}
		\item $D^n \sim pt$ : Lets fix a point $c \in D^n$ which we can think of as the center of this disk. Define a map $f : D^n \rightarrow \{c\}$ to be the constant map taking everythig in the disk to a point $c$. Define $g: \{c\} \rightarrow D^n$ to be a map which takes $g(c) = c$. Per our claim it remains to define a homotopy map $H: D^n \times [0,1] \rightarrow D^n$ such that $H(x,0) = f \circ g (x)$ and $H(x,1) = \id_X$. If we take the straight line homotopy
		\[
		H(x,t) = (1-t)c + tx
		\]
		then we immediately satisfy those conditions. (I believe this generalizes further to the fact that any path-connected space is contractible to a point). 
		\item $T \sim S^1$ where $T$ is a solid torus. Intuitively the solid torus can be seen as $D^2 \times S^1\sim pt \times S^1$. What we have can naturally be identified with just $S^1$. 
	\end{enumerate}
	\begin{definition}[Deformation Retraction]
		Let $A \subset X$, we refer to the inclusion as a map $i$, and define $r : X \rightarrow A$. If $i \circ r = \id_A$ then we call this a retraction. If $r \circ i \sim \id_X$ then we call this a deformation retraction. 
	\end{definition}
	Primary example to consider the cylinder $X \times I$. The retraction map $r : B \times I \rightarrow B$ simply drops the interval paramter, then $r \circ i ( b) = r(b \times I) = b$. We can simply define a straight line homotopy and we get the deformation retraction homotopy equivalence. Note here that $B \times I \not \simeq B$ yet $B \times I \sim B$. 	In fact these two maps are equivlent in this case making this a \emph{strong} deformation retraction. With this language we can precisely define what it means for a space to be contractible.
	\begin{definition}
		A space $X$ is called contractible if $X \sim pt$. Alternately whenever there exists a deformation retraction to a point $x_0 \in X$ homotopic to $\id_X$. 
	\end{definition}
	We already saw above an example of the $n$-disk being contractible. In fact if a space is path-connected then it is contractible. Now consider the set $\pi(X,Y) = E(X,Y)/\sim$ where $\sim$ is homotopy equivalence. That is we are stratifying the path space between two topological spaces into a new set where equivalence classes correspond to homotopy equivalent functions. We can check that the properties of an equivlance relation:
	\begin{itemize}
		\item Transitivity corresponds to composition of paths
		\item Symmetry correspond to the reversal of the direction of a path
		\item Reflexivity corresponds to the fact that a path is always homotopic to itself
	\end{itemize}
	Now we give a categorical view of this homotopy equivalence property. Fix a space $Z$, then $\pi(Z, -)$ and $\pi(-,Z)$ define a covariant and contravariant functor respectively from the category of topologicla spaces to the category of sets. That is on objects $X \mapsto \pi(Z,X)$ and $X \mapsto \pi(X,Z)$. Note that for homotopic maps $\phi \sim \phi'$ in $\pi(Z,X)$ that composition with a map $f: X \rightarrow Y$ preserves the homotopy, i.e. $f \circ \phi \sim f \circ \phi '$. Then to $f$ we associate a map $f_*^Z:\pi(Z,X)\rightarrow\pi(Z,Y)$ where $[\phi] \mapsto [f \circ \phi]$. Similarly if we fix $Z$ in the second coordinate we get a map $f^*_Z : \pi(Y,Z) \rightarrow \pi(X,Z)$. If we have a map $h: Z \rightarrow W$, then we ge a collection of maps which behave naturally in the sense that 
	% https://q.uiver.app/#q=WzAsOCxbMCwwLCJcXHBpKFcsWCkiXSxbMiwwLCJcXHBpKFosWCkiXSxbMCwyLCJcXHBpKFcsWSkiXSxbMiwyLCJcXHBpKFosWSkiXSxbNCwwLCJcXHBpKFgsVykiXSxbNiwwLCJcXHBpKFgsWikiXSxbNCwyLCJcXHBpKFksVykiXSxbNiwyLCJcXHBpKFksWikiXSxbMSwzLCJmXypeWiJdLFswLDIsImZfKl5XIl0sWzEsMCwiaF4qX1giLDFdLFszLDIsImheKl9ZIiwxXSxbNyw1LCJmXipfWiIsMl0sWzYsNCwiZl4qX1ciLDJdLFs0LDUsImhfKl5YIiwxXSxbNiw3LCJoXypeWSIsMV1d
	\[\begin{tikzcd}
		{\pi(W,X)} && {\pi(Z,X)} && {\pi(X,W)} && {\pi(X,Z)} \\
		\\
		{\pi(W,Y)} && {\pi(Z,Y)} && {\pi(Y,W)} && {\pi(Y,Z)}
		\arrow["{f_*^W}", from=1-1, to=3-1]
		\arrow["{h^*_X}"{description}, from=1-3, to=1-1]
		\arrow["{f_*^Z}", from=1-3, to=3-3]
		\arrow["{h_*^X}"{description}, from=1-5, to=1-7]
		\arrow["{h^*_Y}"{description}, from=3-3, to=3-1]
		\arrow["{f^*_W}"', from=3-5, to=1-5]
		\arrow["{h_*^Y}"{description}, from=3-5, to=3-7]
		\arrow["{f^*_Z}"', from=3-7, to=1-7]
	\end{tikzcd}\]
	commutes.
	\begin{theorem}
		$X$ is homotopy equivalent to $Y$ if for all $Z$ there exists such an $\alpha^Z : \pi(Z,X)\rightarrow \pi(Z,Y)$ natural with respect to $Z$. 
	\end{theorem}
	\begin{proof}
		If $X$ and $Y$ are homotopy equivalent then setting $\alpha^Z = f_Z^*$ suffices as the homotopy equivalence ensures that this map is bijecive. Now suppose that such an $\alpha^Z$ exists for all $Z$, then the natural gives us the following commutative square:
		% https://q.uiver.app/#q=WzAsNixbMCwwLCJcXHBpKFgsWCkiXSxbMCwyLCJcXHBpKFgsWSkiXSxbMiwwLCJcXHBpKFksWCkiXSxbMiwyLCJcXHBpKFksWSkiXSxbNCwwLCJcXHBpKFgsWCkiXSxbNCwyLCJcXHBpKFgsWSkiXSxbMCwxLCJcXGFscGhhXlgiLDFdLFswLDIsImdeKl9YIiwxXSxbMSwzLCJnXipfWSIsMV0sWzIsMywiXFxhbHBoYV5ZIiwxXSxbNCw1LCJcXGFscGhhXlgiLDFdLFsyLDQsImZeKl9YIiwxXSxbMyw1LCJmXipfWSIsMV0sWzAsMywiXFxhbHBoYV5ZZ14qX1ggPSBnX1leKlxcYWxwaGFeWCIsMV0sWzMsNCwiKFxcYWxwaGFeWCleey0xfWZfWV4qPWZfWF4qKFxcYWxwaGFeWSleey0xfSIsMV1d
\[
			\begin{tikzcd}
				{\pi(X,X)} && {\pi(Y,X)} && {\pi(X,X)} \\
				\\
				{\pi(X,Y)} && {\pi(Y,Y)} && {\pi(X,Y)}
				\arrow["{g^*_X}"{description}, from=1-1, to=1-3]
				\arrow["{\alpha^X}"{description}, from=1-1, to=3-1]
				\arrow["{\alpha^Yg^*_X = g_Y^*\alpha^X}"{description}, from=1-1, to=3-3]
				\arrow["{f^*_X}"{description}, from=1-3, to=1-5]
				\arrow["{\alpha^Y}"{description}, from=1-3, to=3-3]
				\arrow["{\alpha^X}"{description}, from=1-5, to=3-5]
				\arrow["{g^*_Y}"{description}, from=3-1, to=3-3]
				\arrow["{(\alpha^X)^{-1}f_Y^*=f_X^*(\alpha^Y)^{-1}}"{description}, from=3-3, to=1-5]
				\arrow["{f^*_Y}"{description}, from=3-3, to=3-5]
			\end{tikzcd}
\]
Let $\alpha^X(\id_X)$ denote the homotopy equivalence class of the map $f: X\rightarrow Y$ and $(\alpha^Y)^{-1}$ denote the homotopy equivalence class of $g : Y \rightarrow X$. Then,
\begin{align}
	[\id_X] = \alpha^Yg_X^*(\id_Y) &= g_Y^*\alpha^X(id_Y) \tag{\text{Commutativity}} \\
	&= g_Y^*([f]) \\
	&= [f \circ g]
\end{align}
and 
\begin{align}
	[\id_Y] = (\alpha^X)^{-1}f_Y^*(\id_X) &= f_X^* (\alpha^Y)^{-1}(\id_X) \\
	&= f_X^*([g]) \\
	&= [g \circ f]
\end{align}
The equivalence of these equivalence classes implies directly that $f \circ g\sim\id_Y$ and $g \circ f \sim \id_X$ implying that $X$ is homotopy equivalent to $Y$. 
	\end{proof}
	\section{Homotopy Equivalence of Maps}
	Consider that any map $f:X\rightarrow Y$ factors as follows 
	% https://q.uiver.app/#q=WzAsMyxbMCwwLCJYIl0sWzIsMCwiWFxcdGltZXMgWSJdLFs0LDAsIlkiXSxbMCwxLCJpIiwxLHsic3R5bGUiOnsidGFpbCI6eyJuYW1lIjoiaG9vayIsInNpZGUiOiJ0b3AifX19XSxbMSwyLCJcXHBpIiwxXSxbMCwyLCJmID0gXFxwaSBcXGNpcmMgaSIsMSx7ImN1cnZlIjotM31dXQ==
	\[\begin{tikzcd}
		X && {X\times Y} && Y
		\arrow["i"{description}, hook, from=1-1, to=1-3]
		\arrow["{f = \pi \circ i}"{description}, curve={height=-18pt}, from=1-1, to=1-5]
		\arrow["\pi"{description}, from=1-3, to=1-5]
	\end{tikzcd}\]
	\begin{definition}[Homotopy Equivalence]
		We say $\tilde{f}:\tilde{X}\rightarrow\tilde{Y}$ is homotopy equivalent to $f: X\rightarrow Y$ if there exists maps $\phi: \tilde{X}\rightarrow Y$, $\psi : \tilde{Y}\rightarrow Y$ such that 
		$f \circ \phi \sim \psi \circ \tilde{f}$. 
		% https://q.uiver.app/#q=WzAsNCxbMCwwLCJcXHRpbGRlIFgiXSxbMiwwLCJcXHRpbGRlIFkiXSxbMCwyLCJYIl0sWzIsMiwiWSJdLFswLDEsIlxcdGlsZGUgZiIsMV0sWzAsMiwiXFxwaGkiLDFdLFsxLDMsIlxccHNpICIsMV0sWzIsMywiZiIsMV1d
		\[\begin{tikzcd}
			{\tilde X} && {\tilde Y} \\
			\\
			X && Y
			\arrow["{\tilde f}"{description}, from=1-1, to=1-3]
			\arrow["\phi"{description}, from=1-1, to=3-1]
			\arrow["{\psi }"{description}, from=1-3, to=3-3]
			\arrow["f"{description}, from=3-1, to=3-3]
		\end{tikzcd}\]
	\end{definition}
	Suppose that we have a map $f:X\rightarrow Y$, take $\tilde Y = cyl(f)$ and note that $r: cyl(f)\rightarrow Y$ is a strong deformation retraction then our square becomes
	% https://q.uiver.app/#q=WzAsNCxbMCwwLCJYIl0sWzIsMCwiY3lsKGYpIl0sWzAsMiwiWCJdLFsyLDIsIlkiXSxbMCwxLCJpIiwxLHsic3R5bGUiOnsidGFpbCI6eyJuYW1lIjoiaG9vayIsInNpZGUiOiJ0b3AifX19XSxbMCwyLCJcXGlkX1giLDFdLFsxLDMsInIiLDFdLFsyLDMsImYiLDFdXQ==
	\[\begin{tikzcd}
		X && {cyl(f)} \\
		\\
		X && Y
		\arrow["i"{description}, hook, from=1-1, to=1-3]
		\arrow["{\id_X}"{description}, from=1-1, to=3-1]
		\arrow["r"{description}, from=1-3, to=3-3]
		\arrow["f"{description}, from=3-1, to=3-3]
	\end{tikzcd}\]
	If we can show that $f \circ \id_X \sim r \circ i$ then we can construct a homotopy equivalence for any map $f$. Note that $f$ is exactly $r \circ i$ and $\id_X \circ f = f$ and of course a function is homotopic to itself and so we are done. Though we are not forced into using the mapping cylinder, as long as we can factor the map as an inclusion/projection through the object of choice then we can construct the homotopy equivalence. In such a case we ask that given some projection $\pi : E \rightarrow B$ that lives in this homotopy equivalence square, could we lift paths from $B$ to $E$. Suppose that we have a homotopy $f_t : Z \times I \rightarrow B$, then we get a natural inclusion from $Z \times \{0\}$ into $Z \times I$ which we pair with a "lift". This lift is a map $\tilde f_0  : Z \times \{0\} \rightarrow E$ such that $\pi \circ f_0 = \tilde f_0$. Then we are furnished with a "covering homotopy", a family of maps $\tilde f_t$ such that $Z \times I \rightarrow B$ where $\pi \circ f_t = \tilde f_t$. 
	\begin{definition}
		A map $\pi : E \rightarrow B$ is said to have the Covering Homotopy Property (CHP) if we get the square
		% https://q.uiver.app/#q=WzAsNCxbMCwwLCJaIFxcdGltZXMgMCJdLFsyLDAsIkUiXSxbMiwyLCJCIl0sWzAsMiwiWiBcXHRpbWVzIEkiXSxbMSwyLCJcXHBpIiwxXSxbMywyLCJmX3QiLDFdLFswLDMsImkiLDEseyJzdHlsZSI6eyJ0YWlsIjp7Im5hbWUiOiJob29rIiwic2lkZSI6InRvcCJ9fX1dLFswLDEsIlxcdGlsZGUgZl8wIiwxXSxbMywxLCIoXFx0aWxkZSBmX3QpIiwxXV0=
		\[\begin{tikzcd}
			{Z \times 0} && E \\
			\\
			{Z \times I} && B
			\arrow["{\tilde f_0}"{description}, from=1-1, to=1-3]
			\arrow["i"{description}, hook, from=1-1, to=3-1]
			\arrow["\pi"{description}, from=1-3, to=3-3]
			\arrow["{(\tilde f_t)}"{description}, from=3-1, to=1-3]
			\arrow["{f_t}"{description}, from=3-1, to=3-3]
		\end{tikzcd}\]
		More precisely if for any homotopy $f_t$ and a lift $\tilde f_0$ there exists a covering homotopy $(\tilde f_t)$. 
	\end{definition}
	Naturally, as we are in a homotopy theory class, we ask ourselves whether paths in $B$ can be lifted to paths in $E$. The answer is - almost. If $B$ is a path-connected space then the fibers of the projection $\pi^{-1}(x_i)$ are homotopy equivalent to one another. 
	\begin{theorem}
		If a map $\pi : E \rightarrow B$ over a path connected space $B$ has the covering homotopy property then the fibers of the projection $\pi$ are homotopy equivalent. 
	\end{theorem}
	\begin{proof}
		Let $x_1,x_2 \in B$ be points connected with a path $\gamma : [0,1] \rightarrow B$. Denote $F_i = \pi^{-1}(x_i)$ to be the fibers of the projection over elements of $B$. Observe that naturally $F_0 \hookrightarrow E$ and take $Z = F_0$ treating this inclusion as a lift over the homotopy $f: F_0 \times [0,1] \rightarrow B$ where $f = \gamma \circ pr_2$ that is we project to the interval and utilize the path we have defined earlier. The CHP furnishes us with a family of $\tilde f_t : F_0 \rightarrow F_t$ to the fibers of $F_t$ over $\gamma(t)$, note that the choice to map to the fibers over points of the path is natural because we need $\pi \circ \tilde f_t = f(z,t)$. With this we get a map $f_1 : F_0 \rightarrow F_1$. Taking $Z = F_1$ we obtain a family of maps $\tilde g_t : F_1 \rightarrow F_{1-t}$ in particular we will pick out a map $g_1 : F_1 \rightarrow F_0$. To show that $F_1$ and $F_0$ are homotopy equivalent we will show that $f_1$ and $g_1$ are homotopy inverse i.e. $g_1 \circ f_1 \sim \id_{F_1}$ and $f_1 \circ g_1 \sim \id_{F_0}$. We will utilize the CHP property of $\pi$ again by take $Z = F_0 \times [-1,1]$ with a family of maps $(h_t) : Z \rightarrow E$ such that 
		\[h_t(z) = \begin{cases}
			f_{1+t}(z) &, t \leq 0 \\
			g_1 \circ f_1 &,t > 0
		\end{cases}.\]
		\begin{remark}
			Intuitively we are thinking of this as connecting the lifts of the forward path with the reverse path stringing this together into one homotopy path. Then we will choose another parameters that will control the shrinking of the path in $B$ and modify it's corresponding family of paths in $E$.
		\end{remark}
		Observe that with this definition $h_{-1}(z) = f_0(z) = \id_{F_0}$ and $h_1(z) =g_1 \circ f_1 $
		Consider a new homotopy $\delta_u : [-1,1] \times [0,1] \rightarrow B$ such that 
		\[
		\delta_u(t) =\begin{cases}
			\gamma(u(1+t)) &, t\leq 0 \\
			\gamma(u(1-t)) &, t > 0
		\end{cases}.
		\]
		Observe that $\delta_0(t) = \gamma(0) = x_0$ for all $t \in [-1,1]$ effectively acting as the costant path. Now a homotopy $F_0 \times [-1,1] \xrightarrow{pr_2} [-1,1] \xrightarrow{\delta_u} B$, via the CHP, can be lifted to a family of maps parametrized by $t$ and $u$, $h_{t,u} : F_0 \times [-1,1] \times [0,1] \rightarrow E$ where $h_{t}(1) = \delta_1(t)$ for all $t \in [-1,1]$ is equivalent to the path $h_t$. Now consider the path traced over segments $\{-1\} \times [0,1]$, $[-1,1] \times 0$, and $\{1\} \times [0,1]$ with endpoints $(-1,1)$ and $(1,1)$. As the maps
		\[h_{-1,1} = h_{-1} = \id_{F_0} \quad\text{and}\quad h_{1,1} = h_1 = g_1 \circ f_1\]
		we have successfully constructed a homotpy equivalence of $\id_{F_0} \sim g_1 \circ f_1$. 
		\end{proof}
		As we established before this, for any map $f: X \rightarrow Y$ we can always construct a homotopically equivalent map $\tilde f : X \rightarrow Cyl(f)$ but we can actually simplify this by instead taking $E  = X \times E(Y) = \{(x,\gamma) : \gamma(1) = f(x)\}$. Not only does this give us a homotpy equivalence we can further extend this to show that this has the CHP. This establishes the fact that we only need one map to construct a homotopically equivalent inclusion from a total space to $Y$. 
		\begin{remark}
			My understanding of this is that this allows us to study maps with "nicer" properties without losing topological information, up to homotopy equivalence of course.
		\end{remark}
		Observe that a map $Z \rightarrow E$ is defined utilizing a map $\phi: Z \rightarrow X$ and $\psi : Z \times I \rightarrow Y$ where $\psi(z,0) = \phi(f(z))$. Now construct a map $\xi : Z \times [1,2] \rightarrow Y$ where $\xi(z,1) = \psi(z,1)$. Observe that $\psi(z,t) = \tau_z(t)$ giving us a family of paths $\tau_z \in E(Y)$ parametrized by $z \in Z$. Then to each of these paths we are essentailly just adjoining that path $\xi(z,1)$ raning over $[1,1+t]$. Then it is a matter of reparametrizing the interval and we get a homotopy that lifts automatically for every $t$. 
	\section{CW - Complexes}
		CW Complexes have two ways of being presented but before this I will give an intuitive definition. A CW-Complex is space that is constructed by gluing the Disks together so that boundaries of higher dimensional disks correspond to lower dimensional disks of this space. 
		\begin{definition}[Constructive]
			Let $sk_0(X)$ be a discrete union of disks $D_\alpha^0$, i.e. this is a discrete space and it's cells are disks of dimension 0. Define the $n$-skeleton to be 
			\[
			sk_n(X) = sk_{n-1}(X)\bigcup_\alpha D_\alpha^n
			\]
			where the union with the $\alpha$ is signifying that these spaces are joined via an attaching map $\phi_\alpha: \partial D_\alpha^n \rightarrow sk_{n-1}(X)$. Similarly we can express the $n$ skeleton as a quotient
			\[
			sk_{n-1}(X)\bigsqcup_\alpha D_\alpha^n/\sim
			\]
			where $x \in \partial D_\alpha^n$ is identified with $\phi_\alpha(x) \in sk_{n-1}(X)$. 
		\end{definition}
		$sk_n(X)$ is frequently denoted $X^n$, and with this construction we obtain a nested sequence of inclusion $X^0 \subset X^1 \subset \cdots \subset X^n \subset \cdots$. $X$ is their union equipped with a direct limit topology.
		\begin{definition}[Set-Theoretic]
			Take $X$ to be the union of all cells $\overset{\circ}{D}_\alpha^n$ and for each cell we can define the characteristic map 
			\[
			\Phi_\alpha : D_\alpha^n \rightarrow X
			\]
			where $\partial D_\alpha^n$ coincides with the $\phi_\alpha$.
		\end{definition}
		Similar to direct limit topologies, a subset in $X$ is closed, open, if and if its inverse image in all cells $D_\alpha^n$ is closed, open. 
		\begin{definition}[Axiomatic]
			A space $X$ is equipped with a CW-structure if it is partitioned into pairwise disjoint subsets $e_\alpha^n$ provided with the characteristic maps $\Phi_\alpha : D_\alpha^n \rightarrow X$ which supply the homeomorphism $e_\alpha^n \simeq \overset{\circ}{D^n}$ satisfying the following conditions:
			\begin{enumerate}
				\item Closure-Finiteness:  Each $\Phi_\alpha|_{\partial D^n}$ lands in finitely many cells of lower dimensions
				\item Weak Topology A subset in $X$ is closed iff its inverse image under each characteristic map is closed. 
			\end{enumerate}
		\end{definition}
		At this point it's a good idea to give a good example. As the nested inclusions hint, we should be able to construct the $*^\infty$ spaces using skeletons. First we will construct a finite space, $S^2$. Let $X_0 = \{p\}$ be a single point, then take $X_1 = X_0$, and finally we can take $X_2 = D^2$ with a constant attaching map $\phi: D_2 \rightarrow \{p\}$ be the constant map mapping every point to $p$. Of course this means that the map behaves as we expect on the boundary of $D^2$. Note that there multiple ways to construct a CW-complex for $S^2$. 
		\section{Schubert Cells}
		Given a grassmanian $G(n,k)$ we want to exhibit the CW-structure that this space has naturally. Consider $G(10,5)$ an for an arbitrary $n \times k$ matrix whose rows represent a linear subspace $V^k \cap \R^n$ then the reduced row-echelon form of this matrix is of the form 
		\[
		\begin{bmatrix}
			& \note{1} & * & * & 0 & * & 0 & 0 & * & 0 \\
			&  &  &  & \note{1} & * & 0 & 0 & * & 0 \\	
			&  &  &  &  &  & \note{1} & 0 & * & 0 \\
			&  &  &  &  & &  & \note{1} & * & 0 \\
			&  &  &  &  &  &  &  &  & \note{1} \\
		\end{bmatrix}
		\]
		Let $\lambda = \{2, 5, 7,8,10\}$ be the number representing the distance of the first non-zero number of a row from the left edge of the matrix (including the number itself). 	Observe that as we gradually move from right to left on this matrix, we can think of the output as the intersection of this subspace $V^5$ with $E_i = Span\{e_{i + 1}, \dots, e_n\} = \R^{10-i}$. We can record these dimensions in a list 5,5,4,4,3,3,2,1,1,0. This is a signature, which I am calling a Schubert signature but I don't know what it is called, that gives us the Schubert cells. That is we take a collection of linear subspaces of dimension five whose signature is the same. From this signature we can extract a sequence $\lambda$ that we show above, where $\lambda_i$ tells us after which dimensions of $E_i$, the dimension of $V^5 \cap E_i$ will decrease. By removing the free columns, we recover a Young Tableu/Diagram corresponding to a schubert cell. In turn this relationship allows us to count the number of Schubert cells, which are the number of Young Diagrams that we can fit in a $n-k \times k$ matrix, exactly $n \choose k$. Despite the fact that the schubert cells are rather easy to define, the attaching maps present a problem and it will be dealt with when we get to Morse theory in a very difficult theorem that we will, for the time being, accept.
		
		\section{Bruhat Cells}
		Flag manifolds $V^{d_1} \subset V^{d_2} \subset \dots \subset V^{d_r} \subset \mathbb{K}^n$, consist of a sequence of nested subspaces of $\R^n$, that is a choice of $V^{d_i}$ is a choice of subspace in $ G(n,d_i )$. Then we can imagine that each flag of this flag manifold consists of a schubert cell. That is take $d_{ij} = V^{d_i} \cap E_j$ and record for all $j$, and stick this as a column vector in a matrix. For complete flag manifolds the end result is an $n \times n$, what I am calling a rank-matrix, matrix
		\[
		\begin{bmatrix}
			d_{1,1} & \cdots & d_{1,n} \\
			\vdots & & \vdots \\
			d_{n,1} & \cdots & d_{n,n}
		\end{bmatrix} = LPU
		\]
		with an $LPU$-factorization. The location of the pivot positions of this matrix tell us at which $j$ the intersection $V^{d_i} \cap E_j$ is empty, when we have passed its maximum dimension. This fixes some pivot positions related to the structure of the subspace $V^{d_i}$ and note that pivot positions are determined by the permutation matrix $P$. So in reverse, given an $LPU$ factorization we parametrize a family of complete flag manifolds based on permutations of the identity matrix. If we were to consider a simple example 
		\[
		V^1 \subset V^2 \subset \R^3
		\]
		we would get a $3 \times 3$ rank matrix. This tells us that we have the number of Bruhat cells equivalent to the number of purmutations of $P$ which in this case is $3!$.  
		
		\section{Homotopy Theory of Cell Spaces}
		\begin{definition}
			A pair $(X,A)$ is a Borsuk pair if any map from $X$ to $Y$ and a homotopy of it's restriction to $A$ can be extended to a homotopy on $X$. In other words $X \cup (A \times I) \rightarrow Y$ extends to $X \times I \rightarrow Y$. 
		\end{definition}
		This is the definition provided but it's worth decomposing what we actually have. Given a map $f: X \rightarrow Y$ and a homotopy $h : A \times I \rightarrow Y$ such that $h(a,0) = f(0)$ then we can combine these to form a homotopy on $X \cup (A \times I)$. Note that this is really the space $X \times \{0\} \cup A \times I$ which have an overlap at $t = 0$. Then combining $f$ and $h$ our homotopy is $H : X \cup (A \times I) \rightarrow Y$ such that $H(x,0) = f(x)$ and $H(a,t) = h(a,t)$ which guarantees that the map is a homotopy of the restriction $f|_A$. Then $(X,A)$ is a Borsuk pair if we can extend to a homotopy $\hat F : X \times I \rightarrow Y$ such that $\hat{F}|_{X \cup (A \times I)} = H$. 
		
		\begin{theorem}
			A CW-pair $(X,A)$ is Borsuk
		\end{theorem}
		This arugment will be proved using induction on cells. Later we will see that some spaces are homotopy equivalent to CW-complexes which allows us to prove things, up to homotopy, inductively about these spaces that we otherwise couldnt. We can choose to index the cells $n - 1$ so that $X^{n-1}$ when $n = 1$ is the empty set and our condition is vacuously true and all we have to prove is the induction step. 
		\begin{proof}
			We will be doing induction on the $n$-cells of $X$. Indexing each cell by $n -1$ then $X^{n-1} = \emptyset $ for $n = 1$ so our base case is vacuously true. Suppose that this is true for some $n- 1$ and $X^{n-1 } \subsetneq A$ it will suffice to show that we can construct a homotopy $X^n \cup( (A) \times I) \rightarrow Y$ corresponds with the restriction to the $n-1$ cell. More precisely, the map $f:X\rightarrow Y$ furnishes maps $f_n : X^n \times \{0\} \rightarrow Y$ and for $n-1$ we have homotopy $F_{n-1} : X^{n-1} \cup (A \times I ) \rightarrow Y$ which behaves properly with restriction to lower dimensional cells of $X$. Let $\Phi : D^n \rightarrow X$ be the characteristic map and define the function $D : D^n \cup (\partial D^n \times I) \rightarrow Y $ where 
			\[
			F(D^n, 0) := (f_n \circ \Phi)(D^n) \quad\text{and}\quad F(\partial D^n, 0) := (F_{n-1} \circ \Phi)(\partial D^n).
			\]
			The domain space we can descirbe as the empty cup it naturally lives in $D^n \times I$. This means that composed with this retraction we get the desired homtopy for this $n$ space which behaves properly on the restriction to the base space.
		\end{proof}
		\begin{theorem}
			If $(X,A)$ is Borsuk and $A$ is contractible, the the canonical projection $\pi: X \rightarrow X/A$ is a homotopy equivalence. 
		\end{theorem}
		\begin{proof}
			Let $(X,A)$ be a Borsuk pair with a map $f_0 : X \rightarrow X$ where $f_0 = \id_X$. The contractibility of $A$ furnishes a homotopy $h: A \times I \rightarrow X$ where $h_0(A) = f_0(A)$ and $h_1(A) = pt$. Then we can construct a homotopy $H: X \cup (A \times I) \rightarrow X$ where $H_0(X) = f_0(X)$ and $H_t(A) = h_t(A)$. This homotopy can be extended to $F: X \times I \rightarrow X$ such that $F_1(X)$ is a constant map and $F_t(A) \subset A$. This automotically provides a homotopy $g \circ \phi \sim \id_X$. Taking $\hat{f_t} = \pi \circ f_t$ we get a homotopy.   
		\end{proof}
		\section{Cellular Approximation Theorem}
		We say that a map $f:X \rightarrow Y$ between CW-complexes is cellular if $f(X^n) \subset Y^n$. Intuitively a cellular map never raises the dimension of a cell. What we establish with the cellular approximation theorem is that any map between cellular spaces is always homotopic to a cellular map. This means that given arbitrary maps between CW-complexes we can up to homotopy study more well behaved cellular maps instead. This will show up when we talk about fundamental groups of a space. The fundamental group is the algebraic structure that is generated when we quotient a space by the homotopy equivalence generating homotopy classes. When we talk about paths, we saw that many path operations that preserved looked like group operations an in fact the space $\pi_i(X, Y)$ is in fact a group. A path between two CW-complexes passes through some $n$ cells and for $n \geq 2$ the path can be pushed, i.e. is homotopic, to a cell of $n-1$ dimension. This means that the fundamental group of CW-complexes will, up to homotopy, involve one the 1 and 2-skeletons of the CW complexes. 
		\begin{theorem}
			Any map $f:X \rightarrow Y$ between cellular spaces is homotopic to a cellular map. Moreover, if $f$ is already cellular on a CW-subcomplex $A \subset X$, the the homotopy can be made relative to $A$. 
		\end{theorem}
		The proof of this theorem will rely on the following non-trivial lemma. 
		\begin{lemma}
			Let $Y$ be obtained from $B$ by attaching an $n$-cell, assume that we have the pushout diagram 
			% https://q.uiver.app/#q=WzAsNCxbMCwwLCJcXHBhcnRpYWwgRF5uIl0sWzIsMCwiQiJdLFsyLDIsIkUiXSxbMCwyLCJEXm4iXSxbMywyLCJcXFBoaSJdLFswLDFdLFswLDNdLFsxLDJdXQ==
			\[\begin{tikzcd}
				{\partial D^n} && B \\
				\\
				{D^n} && E
				\arrow[from=1-1, to=1-3]
				\arrow[from=1-1, to=3-1]
				\arrow[from=1-3, to=3-3]
				\arrow["\Phi", from=3-1, to=3-3]
			\end{tikzcd}\]
			
			Then any map $f: (\partial D^m, D^m) \rightarrow (B,Y)$ with $m < n$ is homotopic relative to $\partial D^m$ to a map $g$ such that $g(D^m) \subset Y^n$. 
		\end{lemma}
		In a sense this lemma says that if our map is not cellular on a single cell, then we can push it off of the large cell onto a boundary homotopically relative to the boundary. I won't be proving this lemma because the proof is not provided in the text and I am lazy. The idea is that we can always find a free point on the disk, a result of the application of I think Lebesgue's Covering Lemma, that allows us to construct a deformation retract of the disk onto the boundary taking the image of our map with it (?).  Instead I will attempt to provide a proof of the cellular approximation theorem.
		\begin{proof}
			Let $f: X\rightarrow Y$ be a map of CW-complexes and I want to show that I have a map $g_n$ such that $g_n$ is cellular on a cells $X^n$ excluding a subcomplex $A$ and that there exists a homotopy $H_{n}$ of $g_{n-1} \sim_{X^{n-1}} g_n$. As usualy we index so that we start with the empty cell and so our base case is vacuously true. Then for our inductive step we can assume we have a map $g_{n-1}$ cellular on $X^{n-1} \setminus A$ and a homotopy $H_{n-1}$ and we will use these maps to construct $g_n$ and $H_n$. Now consider cells $e_\alpha^n \in X^n$, if for all such cells $g_{n-1}(e_\alpha^n) \subset Y^n$ then we can extend $H_{n-1}$ by setting $g_n = g_{n-1}$ and using the constant homotopy. Otherwise consider that the cells under $g_{n-1}$ land in some subcomplex of higher dimension $N$ where $Y^n \subseteq Y^N \subset Y$ for $N > n$. Take maximal $N$ and choose cell $e_\alpha^n$ such that $g_{n-1}(e_\alpha^n) \cap Y^N \neq \emptyset$ then we get a diagram 
			% https://q.uiver.app/#q=WzAsNixbMCwwLCJcXHBhcnRpYWwgRF5uIl0sWzIsMCwiWF57bi0xfSJdLFsyLDIsIlhebiJdLFswLDIsIkRebiJdLFs0LDAsIllee24tMX0iXSxbNCwyLCJZXk4iXSxbMywyLCJcXFBoaSJdLFswLDEsIlxcUGhpIl0sWzAsMywiIiwxLHsic3R5bGUiOnsidGFpbCI6eyJuYW1lIjoiaG9vayIsInNpZGUiOiJ0b3AifX19XSxbMSwyLCIiLDEseyJzdHlsZSI6eyJ0YWlsIjp7Im5hbWUiOiJob29rIiwic2lkZSI6InRvcCJ9fX1dLFsyLDUsImdfXFxhbHBoYSIsMV0sWzEsNCwiZ197bi0xfSIsMV0sWzQsNSwiIiwxLHsic3R5bGUiOnsidGFpbCI6eyJuYW1lIjoiaG9vayIsInNpZGUiOiJ0b3AifX19XV0=
			\[\begin{tikzcd}
				{\partial e_\alpha^n} && {X^{n-1}} && {Y^{n-1}} \\
				\\
				{e_\alpha^n} && {X^n} && {Y^N}
				\arrow["\Phi", from=1-1, to=1-3]
				\arrow[hook, from=1-1, to=3-1]
				\arrow["{g_{n-1}}"{description}, from=1-3, to=1-5]
				\arrow[hook, from=1-3, to=3-3]
				\arrow[hook, from=1-5, to=3-5]
				\arrow["\Phi", from=3-1, to=3-3]
				\arrow["{g_\alpha}"{description}, from=3-3, to=3-5]
			\end{tikzcd}\]
		\end{proof}
		supplying us with a map $f_{n,\alpha} = \Phi \circ g_{n-1} : (\partial e_\alpha^n, e_\alpha^n) \rightarrow (Y^{n-1}, Y^N)$. Then lemma 7.2 says there exist a homotopy relative to the boundary of the cell to a map $g_{n,\alpha}$ that is cellular on any disk of dimension $m < N$. Repeating this process and gluing each homotopy results in a final homotopy $H_{n,\alpha} : g_{n-1} \sim_{\partial e_\alpha^n} f_{n,\alpha} : e_\alpha^n \rightarrow Y$. This is actually a homotopy on a subcomplex of $X^n$, together which form a Borsuk pair meaning that we can extend this homotopy to $\tilde{H}_{n}$ that agrees with the homotopy on the subcomplex. Setting $g_n = H_n(-,1)$ we are finished. 
		\section{$N$-connected Cell Spaces} 
		\begin{definition}
			A space is $n$ connected if for all $k \leq n$, any map $S^k \rightarrow X$ can be extended to a map $D^{k+1} \rightarrow X$. 
		\end{definition}
		Lets consider $X = \R^2 \setminus \{(0,0)\}$ then a map $S^0 \rightarrow X$ can be extended to $D^1 \rightarrow X$ which is the equivalent of finding a path between two points and that $X$ is 0-connected, this is telling us that being 0-connected is equivalent to being path-connected. But if we were to take $S^1 \rightarrow X$ where $S^1$ is centered at the origin, we cannot extend this to a map $D^2 \rightarrow X$ because we are missing a point and so our space is not 1-connected. Note here that 1-connectedness is equivalent to being path connected. The idea of $n$-connectedness is generalizing this relationship to higher dimensions. 
		\begin{theorem}
			An $n$-connected CW-complex is hootopy equivalent to a CW-complex whose $n$-skeleton consists of one point. 
		\end{theorem}
		Let us consider the CW-complex $S^2$. The CW-complex construction of $S^2$ consists of exactly two points, a point on the 1-skeleton and a point on the 2-skeleton which corresponds to the 1 and 2 connectedness of $S^2$. We will utilize the fact that quotienting $Y$ by a subcomplex $A$ preserves the CW structure and so we will construct $A$ so that $[A]$ is the $n$-skeleton of $Y$. Then utilizing a corollary of the Borsuk property, $Y \sim Y/A$ and the $n$-skeleton will be a point $[A]$.
		\begin{proof}
			content...
		\end{proof}
		
		
			\end{document}